\documentclass[preprint,12pt,authoryear]{elsarticle}

%% Use the option review to obtain double line spacing
%% \documentclass[authoryear,preprint,review,12pt]{elsarticle}

\usepackage{amssymb}
\linespread{1.2}
\usepackage{amsmath}
\usepackage{lineno}
\usepackage{graphicx}
\usepackage[colorlinks=true, allcolors=blue]{hyperref}
\usepackage[a4paper, total={6in, 9in}]{geometry}
% \linenumbers
\journal{Estuaries and Coasts}

\begin{document}

\section*{Seepage models}

\subsection*{Background}



 Seepage is calculated each day from a control-volume balance over the basin:
\[
S \approx Q_r - \frac{dV}{dt},
\]
after filtering overwash days and treating evaporation as small relative to $Q_r$ during the analyzed closures. Preliminary analysis indicates seepage is substantial, accounting for roughly 20--60\% of total discharge.

We assume Darcy's law to model flow through the berm, considering the berm as a porous body with cross-sectional flow area $A_b$ and flow length $L_b$ from estuary to ocean. The seepage rate $S$ $[L^3T^{-1}]$ scales with the head difference between estuary and ocean,  $\Delta h = h - h_{\rm ocn}$:
\[
S = K\,A_b\,\frac{\Delta h}{L_b},
\]
where $K$ $[LT^{-1}]$ is hydraulic conductivity. Considering that $A_b$ and $L_b$ may vary depending on estuary water level, we define an {integrated berm conductance} $\mathcal{C}$:
\[
\mathcal{C}(\cdot) \;\equiv\; K\,\frac{A_b}{L_b}\qquad [L^2T^{-1}],
\qquad
S \;=\; \mathcal{C}(\cdot)\,\Delta h,
\]
where the conductance  $\mathcal{C}(\cdot)$ may vary with estuary level $h$ or with $\Delta h$.

\subsection*{Transmissivity connection}

The conductance $\mathcal{C}(\cdot)$ is closely related/analogous to aquifer transmissivity, which is a measure of an aquifer's ability to transmit groundwater. Transmissivity is  defined as the product of the aquifer's hydraulic conductivity $K$ and its saturated thickness $h_b$, that is, 
 \(T=K h_b \). 

 If $h_b$ is the vertical dimension of the seepage area and the alongshore width of the active flow section is $W_b$, then $A_b \approx h_b\,W_b$ and
\[
\mathcal{C}(\cdot) \;=\; K\,\frac{A_b}{L_b}
   \;=\; (K\,h_b)\,\frac{W_b}{L_b}
   \;=\; T\, \cdot \frac{W_b}{L_b} ,
\]
where $T=K\,b$ $[L^2T^{-1}]$ is the transmissivity, and if $W_b/L_b$ is nearly constant, $\mathcal{C} \propto T$. 

We parameterize
\[
\mathcal{C}(\cdot) \;=\; C_{\rm ref}\, f(\cdot),
\]
where $C_{\rm ref}$ $[L^2T^{-1}]$ is a baseline (reference) conductance and the dimensionless factor $f(\cdot)$ [-] parameterizes how the seepage geometry ($A_b/L_b$) varies depending on estuary water level or $\Delta h$.


\subsection*{Three candidate models}

In the first, linear model, $f = 1$.
The shifted-power models, we consider that  cross secrtional seepage area  $A_b$ and path length $L_b$ may vary depending on water level. 

\paragraph{(1) Constant conductance (linear).}
\[
f(\cdot)=1
\quad\Rightarrow\quad
S \;=\; C_L\,\Delta h,
\]
with $C_L$ $[L^2T^{-1}]$ constant.

\noindent\textit{Advantage:} Simple, convenient for operations. \\
\textit{Assumption:} The effective seepage cross-sectional area /path is approximately constant over the analyzed range of $\Delta h$.

\paragraph{(2) Visitor-level response (shifted--power on $h$)}

\[
\mathcal{C}(\cdot)= C_h\,f_h(h),
\qquad
f_h(h) \;=\; z_{h}^{\,b_{h}},
\]
\[
z_{h}\;=\;\frac{h_{V}+\theta_{h}}{\,m_{h}\,},\qquad
m_{h}\;=\;\operatorname{median}(h_{V})+\theta_{h},
\]
\[
\Rightarrow\qquad
S \;=\; C_h\,\Delta h\, z_{h}^{\,b_{h}}.
\]
Here $C_h$ $[L^2T^{-1}]$ is a reference conductance, 
$\theta_h$ is an offset (threshold) chosen so $h+\theta_h>0$; $b_h\ge 0$ governs the sensitivity to $h$.

The offset $\theta_h$ represents a minimum water level for flow (that is, a minimum estuary water level below which $A_b$ goes to zero, i.e, when $h + {\theta_h} = 0$ ).  
The $A_b$ is anticipated to increase with rising water level, and the $L_b$ is anticipated to decrease, so $b_h> 0 $ is expected.

Note that, averaged over time, $h- \Delta h \approx 1$  m, because mean sea level is 1 m NAVD at this location.   



\paragraph{(3) Head-difference response (shifted--power on $\Delta h$).}
\[
\mathcal{C}(\cdot)= C_\Delta\,f_\Delta(\Delta h),
\qquad
f_\Delta(\Delta h) \;=\; z_{\Delta}^{\,b_{\Delta}},
\]
\[
z_{\Delta}\;=\;\frac{\Delta h+\theta_{\Delta}}{\,m_{\Delta}\,},\qquad 
m_{\Delta}\;=\;\operatorname{median}(\Delta h)+\theta_{\Delta},
\]
\[
\Rightarrow\qquad
S \;=\; C_\Delta\,\Delta h\, z_{\Delta}^{\,b_{\Delta}}.
\]
Here $C_\Delta$ $[L^2T^{-1}]$ is the reference conductance; $\theta_\Delta$ ensures $\Delta h+\theta_\Delta>0$; $b_\Delta\ge 0$ is estimated separately from $b_h$.

\subsection*{Bounds in estimating model parameters}

We do \emph{not} force $A_b=0$ at $h=0$ (the visitor gauge datum is NAVD88, not estuary depth; $h=0$ $\approx$ mean low tide can still have nonzero wetted thickness). Practically, we bound offsets globally so the powered bases remain positive and physically reasonable (e.g., $\theta \ge -\min(h)$).

\subsection*{Notes}

\begin{itemize}
\item Use $L_b$ consistently for flow length.
\item Use $\mathcal{C}(\cdot)$ for the conductance \emph{function} and $C_{\square}\in\{C_L,C_h,C_\Delta\}$ for the dimensionful reference conductances $[L^2T^{-1}]$.
\item $z$ factors are dimensionless (median $+$ offset), so $C_{\square}$ retains physical units.
\end{itemize}


\end{document}


Since river discharge, estuary height and hypsometry (stage-volume relationship)
data are available, seepage can be estimated by mass balance.
ICEs function as tidal systems when there is an ocean connection, and as salt-stratified
lakes during closures, which prevents vertical mixing and can lead to hypoxia.
An integrated hydraulic conductance can be obtained by regression of seepage on
the head difference between estuary and ocean.
Bar-built estuaries with relatively small (cross sectional areab100 m2) and shallow
tidal inlets are wide- spread in Mediterranean climates and along wave-exposed coasts.

Small, sandy coastal lagoons are common along the U.S. Pacific Coast, provide
valuable habitat for salmonids, and are being squeezed by ongoing sea level rise and
coastal development.